%%% ============================================================
%%% A structural conjecture for the kissing number in dimension
%%% six, with an explicit candidate magic function on Gamma_0(3)
%%%
%%% B.R. Sanders, M. Gish (2026)
%%% Target venue: Journal of Combinatorial Theory, Series A
%%%               (alternative: Algebraic Combinatorics / DCG)
%%% MSC: 52C17, 11F11; secondary 11H31, 11F20
%%%
%%% Draft of 2026-06-10.  TODO markers (red) flag citations or
%%% numerical values that must be confirmed before submission.
%%% ============================================================

\documentclass[11pt,reqno]{amsart}

\usepackage[utf8]{inputenc}
\usepackage[T1]{fontenc}
\usepackage{amsmath,amssymb,amsthm,mathtools}
\usepackage[margin=1in]{geometry}
\usepackage{booktabs}
\usepackage{array}
\usepackage{xcolor}
\usepackage[colorlinks=true,linkcolor=blue,citecolor=blue,urlcolor=blue]{hyperref}
\usepackage{microtype}
\usepackage{graphicx}

%% -------- theorem environments --------
\theoremstyle{plain}
\newtheorem{theorem}{Theorem}[section]
\newtheorem{lemma}[theorem]{Lemma}
\newtheorem{proposition}[theorem]{Proposition}
\newtheorem{corollary}[theorem]{Corollary}
\newtheorem{conjecture}[theorem]{Conjecture}
\newtheorem{openproblem}[theorem]{Open Problem}

\theoremstyle{definition}
\newtheorem{definition}[theorem]{Definition}
\newtheorem{observation}[theorem]{Observation}

\theoremstyle{remark}
\newtheorem{remark}[theorem]{Remark}

%% -------- macros --------
\newcommand{\R}{\mathbb{R}}
\newcommand{\Z}{\mathbb{Z}}
\newcommand{\C}{\mathbb{C}}
\newcommand{\HH}{\mathbb{H}}
\newcommand{\eps}{\varepsilon}
\newcommand{\Dthree}{\Delta_{3}}
\newcommand{\psip}{\psi_{+}}
\newcommand{\Ip}{I_{+}}
\newcommand{\Imm}{I_{-}}
\newcommand{\fhat}[1]{\widehat{#1}}
\newcommand{\TODO}[1]{\textcolor{red}{\textbf{[TODO: #1]}}}
\DeclareMathOperator{\res}{res}
\DeclareMathOperator{\ord}{ord}

%% -------- title matter --------
\title[A conjecture for the kissing number in dimension six]%
{A structural conjecture for the kissing number in dimension six,
with an explicit candidate magic function on
\texorpdfstring{$\Gamma_0(3)$}{Gamma0(3)}}

\author{Brayden R.\ Sanders}
\address{7Site LLC, Hot Springs, Arkansas, USA}
\email{brayden@7site.co}

\author{M.\ Gish}
\address{Independent researcher, Hot Springs, Arkansas, USA}
\email{monica.gish1992@gmail.com}

\date{June 10, 2026}

\keywords{kissing number, $E_6$ root system, linear programming bound,
modular forms, eta quotient, Atkin--Lehner involution, Fricke involution,
Viazovska's method}

\subjclass[2020]{Primary 52C17, 11F11; Secondary 11H31, 11F20}

\begin{document}

\begin{abstract}
Let $K(\R^6)$ denote the kissing number of six-dimensional Euclidean
space.  The minimal vectors of the root lattice $E_6$ give
$K(\R^6)\ge 72$, while the best published upper bounds, obtained by
semidefinite programming, are $77$--$78$.  We conjecture
$K(\R^6)=72$ and, in the spirit of Viazovska's solution of the
sphere-packing problem in dimension $8$, we propose an explicit
candidate ``magic function'': a linear combination of Laplace-type
transforms of two weight-$6$ modular objects on $\Gamma_0(3)$, namely
the unique normalized cusp form
$\Dthree=\eta(\tau)^6\eta(3\tau)^6$ (the newform \textsf{3.6.a.a},
Atkin--Lehner sign $-1$) and the meromorphic Fricke-even quotient
$\psip=(E_6(\tau)^2-729\,E_6(3\tau)^2)/\Dthree$.
We prove or computationally verify the properties of these building
blocks that the construction requires: the Hecke eigenform structure of
$\Dthree$ with a perfect match against the LMFDB; the Atkin--Lehner
eigenvalue $-1$, confirmed four independent ways; consistency with the
Ramanujan--Petersson bound at all primes $p\le 97$; the simple poles of
$\psip$ at both cusps with leading Laurent datum
$-728\,q^{-1}-5376$, where $-728=-2^3\cdot7\cdot13$ and
$-5376=-2^8\cdot3\cdot7$; the divisibility of every Laurent coefficient
of $\psip$ by $56$; a forced zero of $\psip$ at the Fricke fixed point
$i/\sqrt3$; and a $30$-digit numerical profile of the cusp-form
transform.  The construction is the natural level-$3$ analogue of the
level-$1$ machinery used in dimension $8$.  We emphasize that this
paper does \emph{not} prove $K(\R^6)=72$: the analytic continuation of
the transform of $\psip$ and the linear-programming positivity
conditions remain open, and we state the remaining problems precisely.
\end{abstract}

\maketitle

%% ============================================================
\section{Introduction}\label{sec:intro}
%% ============================================================

\subsection{The kissing number problem}

The \emph{kissing number} $K(\R^n)$ is the maximal number of
nonoverlapping unit balls in $\R^n$ that can simultaneously touch a
central unit ball.  Equivalently, after rescaling, $K(\R^n)$ is the
maximal cardinality of a set of points on a sphere with all pairwise
angular distances at least $60^\circ$.  The problem is classical and,
in most dimensions, open.  The known sharp values are summarized in
Table~\ref{tab:known}.

\begin{table}[ht]
\centering
\begin{tabular}{cclc}
\toprule
$n$ & $K(\R^n)$ & optimal configuration & reference \\
\midrule
$1$ & $2$ & $\pm 1$ & elementary \\
$2$ & $6$ & hexagonal & elementary \\
$3$ & $12$ & icosahedral / fcc, non-unique & \cite{SchuetteWaerden1953} \\
$4$ & $24$ & $D_4$ minimal vectors ($24$-cell) & \cite{Musin2008} \\
$8$ & $240$ & $E_8$ minimal vectors & \cite{Levenshtein1979,OdlyzkoSloane1979} \\
$24$ & $196560$ & Leech minimal vectors & \cite{Levenshtein1979,OdlyzkoSloane1979} \\
\bottomrule
\end{tabular}
\medskip
\caption{Dimensions in which the kissing number is known.  Uniqueness
of the optimal configurations in dimensions $8$ and $24$ is due to
Bannai and Sloane \cite{BannaiSloane1981}.}
\label{tab:known}
\end{table}

The sharp values in dimensions $8$ and $24$ were obtained in 1979 by
Levenshtein \cite{Levenshtein1979} and, independently, Odlyzko and
Sloane \cite{OdlyzkoSloane1979}, through the linear-programming (LP)
method of Delsarte, Goethals and Seidel \cite{DGS1977} for spherical
codes: in those two dimensions a low-degree polynomial certificate
matches the lattice lower bound exactly.  No such certificate has
been found in dimension $6$ (or in any other currently open
dimension): the polynomial LP and its known semidefinite-programming
(SDP) strengthenings plateau strictly above the lattice bound.

\subsection{Dimension six}\label{subsec:dimsix}

The root lattice $E_6$, normalized to minimal norm $2$, has $72$
minimal vectors, so
\[
K(\R^6)\;\ge\;72 .
\]
On the other side, the classical LP bound of \cite{OdlyzkoSloane1979}
gives $K(\R^6)\le 82$
\TODO{confirm the exact Odlyzko--Sloane value for $n=6$},
the three-point SDP bound of Bachoc and Vallentin
\cite{BachocVallentin2008} (computed to high accuracy in
\cite{MittelmannVallentin2010}) gives
\[
K(\R^6)\;\le\;78,
\]
and an exploitation of polynomial symmetry by Machado and Oliveira
\cite{MachadoOliveira2018} reportedly sharpens this to $K(\R^6)\le 77$
\TODO{confirm the exact bound and reference data in
\cite{MachadoOliveira2018}}.  Thus
\[
72\;\le\;K(\R^6)\;\le\;78,
\]
and the gap of $5$--$6$ contact points has resisted all
polynomial-certificate methods.  It is widely expected that the lattice
bound is the truth.  We record this as a formal conjecture.

\begin{conjecture}\label{conj:main}
$K(\R^6)=72$, and the optimal configuration is unique up to isometry,
given by the minimal vectors of the root lattice $E_6$.
\end{conjecture}

Conjecture~\ref{conj:main} is folklore; the contribution of this paper
is an \emph{explicit analytic candidate certificate} for it, together
with proofs and computational verifications of the properties of the
certificate's building blocks, and a precise statement of what remains
open.

\subsection{The modular-form method}

In 2017 Viazovska \cite{Viazovska2017} solved the sphere-packing
problem in dimension $8$ by constructing, from modular forms for
$\mathrm{SL}_2(\Z)$, an explicit radial Schwartz function --- a
``magic function'' in the terminology popularized by Cohn
\cite{Cohn2017} --- that makes the Cohn--Elkies linear programming
bound \cite{CohnElkies2003} sharp at the $E_8$ density.  The method
was immediately extended to dimension $24$
\cite{CKMRV2017}, and was later upgraded to a complete Fourier
interpolation framework establishing universal optimality of $E_8$ and
the Leech lattice \cite{CKMRV2022}, through Euclidean interpolation
certificates rather than spherical polynomials.

The lesson of \cite{Viazovska2017,CKMRV2017,CKMRV2022} is that when
polynomial certificates stall above a lattice bound, the correct
certificates may instead be \emph{transcendental} functions built from
the modular forms attached to the lattice.  For $E_8$ and the Leech
lattice the relevant forms live at level $1$.  The lattice $E_6$ has
discriminant $3$, and its theta series is a weight-$3$ form on
$\Gamma_0(3)$ with character $\chi_{-3}$; the natural home for a
dimension-six certificate is therefore level $3$.  This paper carries
out the level-$3$ bookkeeping and isolates a specific candidate.

\subsection{The candidate and the results of this paper}

Write $q=e^{2\pi i\tau}$ for $\tau$ in the upper half-plane $\HH$, let
\[
\Dthree(\tau)\;=\;\eta(\tau)^6\,\eta(3\tau)^6
\;=\;q\prod_{n\ge1}(1-q^n)^6(1-q^{3n})^6
\;=\;q-6q^2+9q^3+4q^4+6q^5-\cdots
\]
be the unique normalized cusp form of weight $6$ on $\Gamma_0(3)$, and
let
\[
\psip(\tau)\;=\;\frac{E_6(\tau)^2-729\,E_6(3\tau)^2}{\Dthree(\tau)}
\;=\;-\frac{728}{q}-5376+195048\,q+18475520\,q^2+\cdots
\]
be the meromorphic weight-$6$ quotient with Fricke eigenvalue $+1$
constructed in Section~\ref{sec:candidate}.  The candidate magic
function is
\begin{equation}\label{eq:f6}
f_6(x)\;=\;\sin^2\!\Big(\frac{\pi|x|^2}{2}\Big)\,
\Big[\alpha\,\Ip(|x|^2)+\beta\,\Imm(|x|^2)\Big],
\qquad x\in\R^6,
\end{equation}
where $\Imm$ and $\Ip$ are Laplace-type transforms of $\Dthree$ and
$\psip$ against the measure $t^2\,dt$
(Definitions~\ref{def:Iminus} and \ref{def:Iplus}), and the scalars
$\alpha,\beta$ are to be determined by the sharpness conditions
$\fhat{f}_6(\sqrt2)=0$ and $f_6(0)/\fhat{f}_6(0)=72$.

The paper is organized as follows.
Section~\ref{sec:structural} presents three elementary structural
observations that select the value $72$ and the level $3$; these are
motivation, not proof.
Section~\ref{sec:candidate} constructs the candidate
\eqref{eq:f6} and proves the elementary analytic facts about its
ingredients (convergence and positivity of the cusp-form transform,
asymptotics of $\psip$ on the imaginary axis, the action of the Fricke
involution).
Section~\ref{sec:verified} establishes the arithmetic of the two
building blocks: part is proved, part is verified computationally with
exact integer arithmetic and cross-validated against the LMFDB
\cite{LMFDB}; the ancillary verification scripts are described in
Appendix~\ref{app:repro}.  Two facts recorded there appear to be new:
a forced zero of $\psip$ at the Fricke fixed point $i/\sqrt3$
(Proposition~\ref{prop:forcedzero}), and the divisibility of every
Laurent coefficient of $\psip$ by $56$
(Proposition~\ref{prop:fiftysix}).
Section~\ref{sec:gap} states precisely what remains open.
Section~\ref{sec:comparison} compares the construction, item by item,
with the dimension-$8$ machinery.
Appendices~\ref{app:hecke}--\ref{app:repro} contain the data tables
and reproducibility information.

\subsection{What this paper does not claim}\label{subsec:noclaim}

We state this prominently, here and in Section~\ref{sec:gap}.
\emph{This paper does not prove $K(\R^6)=72$.}  The transform $\Ip$ is
defined here only in a regularized form and only for $|x|^2>2$; its
analytic continuation to all radii, the verification of the
linear-programming sign conditions, the determination of
$(\alpha,\beta)$, and the reduction of the kissing bound to the
continued function are open problems, stated as
Open Problems~\ref{op:continuation}--\ref{op:descent}.  What the paper
provides is: (i) a specific, structurally motivated candidate, with
every discrete choice (level, weight, cusp form, Fricke eigenvalues,
numerator) forced by the framework of Section~\ref{sec:candidate};
(ii) complete verification of the arithmetic properties the candidate
needs its building blocks to have; and (iii) a precise reduction of
Conjecture~\ref{conj:main} to a stated analytic problem.

\subsection{Notation and conventions}\label{subsec:notation}

$\HH$ denotes the upper half-plane, $q=e^{2\pi i\tau}$, and
$\eta(\tau)=q^{1/24}\prod_{n\ge1}(1-q^n)$ is the Dedekind eta
function.  The normalized Eisenstein series of weight $6$ is
\[
E_6(\tau)\;=\;1-504\sum_{n\ge1}\sigma_5(n)\,q^n ,
\qquad \sigma_5(n)=\sum_{d\mid n}d^5 ,
\]
a modular form of weight $6$ for $\mathrm{SL}_2(\Z)$.
For $N\ge1$, $\Gamma_0(N)$ is the usual congruence subgroup.  The
weight-$k$ slash operator with the determinant normalization is
$(f|_k\gamma)(\tau)=(\det\gamma)^{k/2}(c\tau+d)^{-k}f(\gamma\tau)$.
The \emph{Fricke involution} on level $N$ is
$W_N=\begin{psmallmatrix}0&-1\\N&0\end{psmallmatrix}$, so that at
$N=3$, $k=6$,
\begin{equation}\label{eq:slashW3}
(f|_6W_3)(\tau)\;=\;3^{-3}\,\tau^{-6}\,f\!\big(-\tfrac1{3\tau}\big).
\end{equation}
$W_3$ normalizes $\Gamma_0(3)$ and $W_3^2=1$ on weight-$6$ forms, so
$M_6(\Gamma_0(3))$ splits into $W_3$-eigenspaces with eigenvalues
$\pm1$; since $3$ is prime, $W_3$ is also the unique nontrivial
Atkin--Lehner involution at level $3$ \cite{AtkinLehner1970}.

The Fourier transform on $\R^n$ is
$\fhat{f}(\xi)=\int_{\R^n}f(x)\,e^{-2\pi i\langle x,\xi\rangle}\,dx$;
for radial $f$ we freely write $f(r)$ for the common value on
$|x|=r$.  Throughout, the kissing configuration is normalized as in
$E_6$: the contact points lie on the sphere of radius $\sqrt2$, and
tangency of the outer balls forces pairwise distances $\ge\sqrt2$,
equivalently pairwise inner products $\le 1$, equivalently
$\cos\theta\le\tfrac12$.

%% ============================================================
\section{Structural observations selecting \texorpdfstring{$72$}{72}
and level \texorpdfstring{$3$}{3}}\label{sec:structural}
%% ============================================================

This section collects three elementary observations which together
single out the value $K=72$ within the admissible window
$[72,78]$ and dictate the modular-form home used in
Section~\ref{sec:candidate}.  They are heuristic: we do not know how
to convert any of them into a bound, and nothing in
Sections~\ref{sec:candidate}--\ref{sec:verified} depends logically on
this section.  We include them because they explain why the candidate
of Section~\ref{sec:candidate} was singled out \emph{before} any
numerics were run, and because the same pattern of reasoning may be
useful in other open dimensions.

\subsection{Smoothness of the known optima and of the dual pair}

\begin{observation}\label{obs:smooth}
Every known sharp kissing number factors over the prime set
$\{2,3,5,7,11,13\}$:
\[
2,\quad 6=2\cdot3,\quad 12=2^2\cdot3,\quad 24=2^3\cdot3,\quad
240=2^4\cdot3\cdot5,\quad 196560=2^4\cdot3^3\cdot5\cdot7\cdot13 .
\]
In dimensions $1$--$4$ the optima are in fact $\{2,3\}$-smooth.  In
the dimension-six window $72\le K\le 78$, the integer
$72=2^3\cdot3^2$ is the \emph{unique} $\{2,3\}$-smooth value
($73$ is prime; $74=2\cdot37$; $75=3\cdot5^2$; $76=2^2\cdot19$;
$77=7\cdot11$; $78=2\cdot3\cdot13$).  Moreover, writing $E_6^*$ for
the dual lattice, whose $54$ minimal vectors realize the kissing
number of the dual configuration, the dual-pair product is again pure
$\{2,3\}$:
\[
K(E_6)\cdot K(E_6^*)\;=\;72\cdot 54\;=\;3888\;=\;2^4\cdot3^5 .
\]
\end{observation}

We make no claim that smoothness is more than a regularity of the
known examples: of the seven integers in the window, three
($75$, $77$, $78$) do factor over $\{2,3,5,7,11,13\}$, so it is the
stronger $\{2,3\}$-smoothness --- matching the discriminant $3$ of
$E_6$ and the level-$3$ structure below --- that isolates $72$.  The
$54$ minimal vectors of $E_6^*$ form the classical configuration of
$27$ antipodal pairs corresponding to the $27$ lines on a smooth cubic
surface, on which the Weyl group $W(E_6)$ acts \cite{ConwaySloane1999}.

\subsection{The Eisenstein structure of \texorpdfstring{$E_6$}{E6}}

\begin{observation}\label{obs:eisenstein}
$E_6$ carries the structure of a rank-$3$ Hermitian lattice over the
Eisenstein integers $\Z[\omega]$, $\omega=e^{2\pi i/3}$
\cite{ConwaySloane1999}.  Scalar multiplication by any unit
$\mu\in\{\pm1,\pm\omega,\pm\omega^2\}$, $\mu\ne1$, is an isometry
without nonzero fixed vectors, so the unit group $\mu_6$ acts freely
on the $72$ minimal vectors:
\[
72\;=\;3\cdot 24\;=\;6\cdot 12 ,
\]
i.e.\ the contact points organize into $24$ orbits under
$\langle\omega\rangle$ and into $12$ regular hexagonal frames under
$\mu_6$.
\end{observation}

The factor $24$ matches the kissing number of $\R^4$, and the
hexagonal frames are the two-dimensional optimum; the multiplicative
structure $72=6\cdot12=2\cdot3\cdot12$ is what one expects of an
extremal configuration in dimension $6=2\cdot3$.  Again, this is
numerology with a precise group-theoretic core, not a proof.

\subsection{Discriminant \texorpdfstring{$3$}{3} forces level
\texorpdfstring{$3$}{3}}

\begin{observation}\label{obs:level}
$\det E_6=3$, and the theta series
$\theta_{E_6}(\tau)=\sum_{v\in E_6}q^{|v|^2/2}$ is a modular form of
weight $3$ on $\Gamma_0(3)$ with character $\chi_{-3}$
\cite{ConwaySloane1999,Ebeling2013}.  Any Viazovska-style certificate
adapted to $E_6$ should therefore be built from modular objects on
$\Gamma_0(3)$, just as the dimension-$8$ certificate is built at level
$1$ because $E_8$ is unimodular.
\end{observation}

Observation~\ref{obs:level} is the operative one: it tells us
\emph{where to look}, and Section~\ref{sec:candidate} shows that at
level $3$, weight $6$, the search space collapses to an essentially
unique candidate.

\begin{remark}[Provenance]\label{rem:provenance}
The heuristics of this section arose in a structural framework
developed by the first author for organizing integer invariants of
discrete structures by their factorizations over the prime alphabet
$\{2,3,5,7,11,13\}$.  No familiarity with that framework is needed:
all statements above are elementary and self-contained, and nothing
later in the paper depends on them.
\end{remark}

%% ============================================================
\section{The candidate construction}\label{sec:candidate}
%% ============================================================

\subsection{The linear-programming framework}\label{subsec:LP}

We first recall the rigorous certificate framework our candidate is
designed to enter.  For points on a sphere the classical statement is
the Delsarte--Goethals--Seidel LP bound \cite{DGS1977,KabatianskyLevenshtein1978};
we will use the following Euclidean form, which descends to the sphere
by Schoenberg's theorem.  Let $\sigma$ denote the normalized uniform
measure on the sphere $\sqrt2\,S^5=\{x\in\R^6:|x|=\sqrt2\}$, and for
continuous $f:\R^6\to\R$ set
\[
c_0(f)\;=\;\int\!\!\int f(u-v)\,d\sigma(u)\,d\sigma(v).
\]

\begin{proposition}[Euclidean kissing certificate]\label{prop:certificate}
Let $f:\R^6\to\R$ be a continuous, integrable, radial function such
that
\begin{enumerate}
\item[(i)] $\fhat{f}(\xi)\ge0$ for all $\xi\in\R^6$, with
$\fhat f\in L^1$;
\item[(ii)] $f(x)\le0$ for all $\sqrt2\le|x|\le2\sqrt2$;
\item[(iii)] $c_0(f)>0$.
\end{enumerate}
Then
\[
K(\R^6)\;\le\;\frac{f(0)}{c_0(f)} .
\]
\end{proposition}

\begin{proof}
Let $u_1,\dots,u_N\in\sqrt2\,S^5$ be a kissing configuration, so
$|u_i-u_j|\ge\sqrt2$ for $i\ne j$; automatically
$|u_i-u_j|\le2\sqrt2$.  By (i) and Fourier inversion, $f$ is a
continuous positive-definite function on $\R^6$, hence its restriction
$h(\langle u,v\rangle/2):=f(u-v)$ is a continuous positive-definite
zonal kernel on $\sqrt2\,S^5$.  By Schoenberg's theorem
\cite{Schoenberg1942}, $h$ has an absolutely and uniformly convergent
ultraspherical expansion $h(t)=\sum_{k\ge0}c_kG_k^{(6)}(t)$ with all
$c_k\ge0$, where $G_k^{(6)}$ are the Gegenbauer polynomials for $S^5$
normalized by $G_k^{(6)}(1)=1$; orthogonality gives
$c_0=\int\!\!\int h\,d\sigma\,d\sigma=c_0(f)$.  By the addition
theorem each kernel $G_k^{(6)}$ is positive definite on the sphere, so
$\sum_{i,j}G_k^{(6)}(t_{ij})\ge0$ where $t_{ij}=\langle
u_i,u_j\rangle/2$.  Hence
\[
N\,f(0)\;\ge\;\sum_{i,j}h(t_{ij})
\;=\;\sum_{k\ge0}c_k\sum_{i,j}G_k^{(6)}(t_{ij})
\;\ge\;c_0\,N^2 ,
\]
where the first inequality uses (ii) to discard the off-diagonal
terms.  Divide by $c_0N$.
\end{proof}

\begin{remark}\label{rem:c0}
By Parseval, $c_0(f)=\int_{\R^6}\fhat f(\xi)\,
|\widehat{\sigma}(\xi)|^2\,d\xi\ge0$ always holds under (i); the
content of (iii) is strictness.  Note that $c_0(f)$ is \emph{not}
$\fhat f(0)$ in general.  A modern Euclidean route to sharp spherical
bounds of this type is the interpolation machinery of
\cite{CKMRV2022}, which works directly with functions on $\R^n$; for
our candidate, the passage from the design conditions
\eqref{eq:CE-design} below to a certificate in the precise sense of
Proposition~\ref{prop:certificate} is part of the open program
(Open Problem~\ref{op:descent}).
\end{remark}

\subsection{Modular building blocks at level
\texorpdfstring{$3$}{3}, weight \texorpdfstring{$6$}{6}}
\label{subsec:blocks}

The space $M_6(\Gamma_0(3))$ has dimension $3$, with
$\dim S_6(\Gamma_0(3))=1$ and a $2$-dimensional Eisenstein subspace
spanned by $E_6(\tau)$ and $E_6(3\tau)$
\cite{DiamondShurman2005,LMFDB}.  The cusp line is spanned by the eta
quotient
\[
\Dthree(\tau)\;=\;\eta(\tau)^6\,\eta(3\tau)^6\;\in\;S_6(\Gamma_0(3)),
\]
which is holomorphic and nonvanishing on $\HH$ (as $\eta$ is), vanishes
to order $1$ at the cusp $\infty$, and hence --- since the valence
formula for the index-$4$ subgroup $\Gamma_0(3)$ makes the total cusp
vanishing equal to $6\cdot4/12=2$ --- vanishes to order exactly $1$ at
the cusp $0$ as well.  That $\Dthree$ is modular of this weight and
level follows from the standard eta-quotient criteria (see
\cite[\S1.4]{Ono2004}); it is the newform \textsf{3.6.a.a} of the
LMFDB, which lists precisely this eta-quotient presentation
\cite{LMFDB}.  Its arithmetic is the subject of
Propositions~\ref{prop:hecke}--\ref{prop:RP} below.

The Fricke involution acts on the Eisenstein pair as follows.

\begin{lemma}\label{lem:W3eisenstein}
With the normalization \eqref{eq:slashW3},
\[
E_6(\tau)\big|_6W_3=27\,E_6(3\tau),
\qquad
E_6(3\tau)\big|_6W_3=\tfrac1{27}\,E_6(\tau).
\]
Consequently the combinations
\[
G_{\pm}(\tau)\;=\;E_6(\tau)\,\pm\,27\,E_6(3\tau)
\]
satisfy $G_\pm|_6W_3=\pm G_\pm$, and
$\{G_+,G_-\}$ is the $W_3$-eigenbasis of the Eisenstein subspace of
$M_6(\Gamma_0(3))$.
\end{lemma}

\begin{proof}
From the level-$1$ functional equation $E_6(-1/z)=z^6E_6(z)$ with
$z=3\tau$,
\[
(E_6|_6W_3)(\tau)=3^{-3}\tau^{-6}E_6\!\big(-\tfrac1{3\tau}\big)
=3^{-3}\tau^{-6}(3\tau)^6E_6(3\tau)=27\,E_6(3\tau),
\]
and with $z=\tau$,
\[
\big(E_6(3\,\cdot)|_6W_3\big)(\tau)
=3^{-3}\tau^{-6}E_6\!\big(-\tfrac1{\tau}\big)
=3^{-3}\tau^{-6}\tau^{6}E_6(\tau)=\tfrac1{27}E_6(\tau).
\qedhere
\]
\end{proof}

The cusp form lies in the $(-1)$-eigenspace:
$\Dthree|_6W_3=-\Dthree$ (Proposition~\ref{prop:AL}; the one-line
symbolic proof via $\eta(-1/\tau)=\sqrt{-i\tau}\,\eta(\tau)$ is given
there).  Thus the full eigenspace decomposition of $M_6(\Gamma_0(3))$
is
\[
M_6(\Gamma_0(3))^{+}=\C\,G_+,
\qquad
M_6(\Gamma_0(3))^{-}=\C\,G_-\oplus\C\,\Dthree .
\]

\subsection{The meromorphic Fricke-even quotient
\texorpdfstring{$\psip$}{psi+}}\label{subsec:psiplus}

Viazovska's construction in dimension $8$ takes Laplace-type
transforms of \emph{weakly holomorphic} (pole at the cusp) modular
objects whose Fricke/inversion eigenvalue matches the parity of the
target Fourier eigenfunction, with the discriminant cusp form in the
denominator supplying the pole \cite{Viazovska2017}.  At level $3$,
weight $6$, the same recipe forces the following object.  Define
\begin{equation}\label{eq:psiplus}
\psip(\tau)
\;=\;\frac{G_+(\tau)\,G_-(\tau)}{\Dthree(\tau)}
\;=\;\frac{E_6(\tau)^2-729\,E_6(3\tau)^2}{\eta(\tau)^6\,\eta(3\tau)^6}.
\end{equation}
This is a meromorphic modular form of weight $12-6=6$ on
$\Gamma_0(3)$, holomorphic on $\HH$ (the denominator is nonvanishing
there), with poles confined to the cusps.  Its $W_3$-eigenvalue is
forced to be $+1$: among the three Eisenstein bilinears available as
numerators, the eigenvalue arithmetic over the $(-1)$-eigenform
denominator gives
\[
\frac{G_+^2}{\Dthree}\mapsto\frac{(+1)^2}{-1}=-1,\qquad
\frac{G_-^2}{\Dthree}\mapsto\frac{(-1)^2}{-1}=-1,\qquad
\frac{G_+G_-}{\Dthree}\mapsto\frac{(+1)(-1)}{-1}=+1 ,
\]
so $G_+G_-$ is the unique choice (up to scale) producing a Fricke-even
quotient.  The precise statement and proof are
Proposition~\ref{prop:psiplus}.  The Laurent expansion at $\infty$
begins
\begin{equation}\label{eq:psilaurent}
\psip(\tau)\;=\;-\frac{728}{q}\;-\;5376\;+\;195048\,q
\;+\;18475520\,q^2\;+\;508641336\,q^3\;+\;\cdots ,
\end{equation}
with leading coefficient
$-728=(1+27)(1-27)=-2^3\cdot7\cdot13$ inherited from the constant
terms of $G_\pm$; the coefficients are tabulated in
Appendix~\ref{app:laurent}.

\subsection{The two transforms}\label{subsec:transforms}

The dimension enters through the Laplace kernel: in $\R^6$ the
Gaussian $e^{-\pi t|x|^2}$ has Fourier transform
$t^{-3}e^{-\pi|\xi|^2/t}$, and the measure $t^{(n-2)/2}\,dt=t^2\,dt$
is the standard pairing for radial functions in dimension $n=6$.

\begin{definition}\label{def:Iminus}
For $r^2\ge0$ set
\[
\Imm(r^2)\;=\;\int_0^\infty \Dthree(it)\,e^{-\pi r^2t}\,t^2\,dt .
\]
\end{definition}

\begin{lemma}\label{lem:Iminus}
$\Imm(r^2)$ converges absolutely for every $r^2\ge0$ and defines a
strictly positive function.  Moreover there is a constant $C$ such
that
\[
0\;<\;\Imm(r^2)\;\le\;C\,e^{-\pi r/3}\qquad(r\ge1),
\]
so the radial function $f_-(x):=\Imm(|x|^2)$ is continuous, strictly
positive, and integrable on $\R^6$, with rapid (exponential-type)
decay.
\end{lemma}

\begin{proof}
On the imaginary axis $\Dthree(it)=e^{-2\pi
t}\prod_{n\ge1}(1-e^{-2\pi nt})^6(1-e^{-6\pi nt})^6\in(0,1)$, so the
integrand is positive and bounded by $e^{-2\pi t}t^2$ as
$t\to\infty$.  As $t\to0^+$, the functional equation
$\Dthree(i/(3t))=27\,t^6\,\Dthree(it)$ --- the imaginary-axis form of
$\Dthree|_6W_3=-\Dthree$, see Remark~\ref{rem:signs} --- gives
$\Dthree(it)=\tfrac1{27}t^{-6}\Dthree\big(\tfrac{i}{3t}\big)
\le \tfrac1{27}t^{-6}e^{-2\pi/(3t)}$,
which vanishes to all orders; convergence and positivity follow.  For
the decay bound, split at $t=1/r$: for $0<t\le1/r$,
\[
\Dthree(it)\le\tfrac1{27}\,t^{-6}e^{-\pi/(3t)}\cdot e^{-\pi/(3t)}
\le \Big(\tfrac1{27}\sup_{s>0}s^{-6}e^{-\pi/(3s)}\Big)\,e^{-\pi r/3},
\]
so that stretch of the integral is $\le C_1e^{-\pi r/3}$; for
$t\ge1/r$, bound $\Dthree(it)\,t^2\le M$ (continuous, vanishing at
both ends) and integrate: $\int_{1/r}^\infty Me^{-\pi
r^2t}\,dt=Me^{-\pi r}/(\pi r^2)$.  Integrability of $f_-$ on $\R^6$
follows from the decay bound.
\end{proof}

The cusp-form transform is, by itself, \emph{positive definite}:

\begin{lemma}\label{lem:Fpositive}
The Fourier transform of $f_-(x)=\Imm(|x|^2)$ on $\R^6$ is
\[
\fhat{f_-}(\xi)\;=\;27\int_0^\infty \Dthree(iu)\,
e^{-3\pi|\xi|^2u}\,u^5\,du\;>\;0
\qquad\text{for all }\xi\in\R^6 .
\]
\end{lemma}

\begin{proof}
By Tonelli (all factors positive) we may exchange the $x$- and
$t$-integrals:
\[
\fhat{f_-}(\xi)=\int_0^\infty\Dthree(it)\,t^2\,
\big(t^{-3}e^{-\pi|\xi|^2/t}\big)\,dt
=\int_0^\infty\Dthree(it)\,t^{-1}e^{-\pi|\xi|^2/t}\,dt .
\]
Substituting $t=1/(3u)$ and applying
$\Dthree(i/(3u))=27u^6\Dthree(iu)$ gives the display; positivity is
clear since $\Dthree(iu)>0$.
\end{proof}

\begin{remark}\label{rem:rescale}
Lemma~\ref{lem:Fpositive} exhibits the structural difference between
level $1$ and level $3$: the Fricke involution intertwines the
$(t^2dt,\,\pi r^2)$-transform not with itself but with a
$(u^5du,\,3\pi r^2)$-transform --- a rescaling by $3$ in the Gaussian
parameter, reflecting the index-$3$ relation between $E_6$ and
$E_6^*$.  At level $1$ ($E_8$, self-dual) the corresponding
computation returns the same transform, which is why the $\pm1$
Fricke eigenspaces there map directly to $\pm1$ Fourier
eigenfunctions.  Handling this rescaling exactly is one of the
genuinely new features of the level-$3$ problem; see
Section~\ref{sec:gap}.
\end{remark}

In particular $f_-$ alone cannot be a certificate: it is positive
everywhere, while condition (ii) of
Proposition~\ref{prop:certificate} demands negativity beyond
$\sqrt2$.  Sign changes must come from the meromorphic component and
the trigonometric prefactor.

For $\psip$ the integral requires care at \emph{both} ends.

\begin{lemma}\label{lem:psiasym}
On the imaginary axis, $\psip(it)$ is real, and
\[
\psip(it)\;=\;-728\,e^{2\pi t}-5376+O(e^{-2\pi t})
\quad(t\to\infty),
\qquad
\psip(it)\;=\;\frac{728}{27}\,t^{-6}e^{2\pi/(3t)}
+\frac{5376}{27}\,t^{-6}+O\big(t^{-6}e^{-2\pi/(3t)}\big)
\quad(t\to0^+).
\]
In particular $\psip(it)\to+\infty$ as $t\to0^+$ and
$\psip(it)\to-\infty$ as $t\to\infty$.
\end{lemma}

\begin{proof}
The expansion at $t\to\infty$ is \eqref{eq:psilaurent} with
$q=e^{-2\pi t}$.  For $t\to0^+$, the imaginary-axis form of
$\psip|_6W_3=\psip$ (Proposition~\ref{prop:psiplus};
Remark~\ref{rem:signs}) reads
$\psip(it)=-\tfrac1{27}t^{-6}\,\psip\big(\tfrac{i}{3t}\big)$; insert
the $t\to\infty$ expansion of $\psip(i/(3t))$.
\end{proof}

\begin{definition}[Regularized transform]\label{def:Iplus}
For $r^2>2$ set
\[
\Ip(r^2)\;=\;\int_1^\infty \psip(it)\,e^{-\pi r^2t}\,t^2\,dt
\;+\;\int_0^1\Big[\psip(it)-\frac{728}{27}t^{-6}e^{2\pi/(3t)}
-\frac{5376}{27}t^{-6}\Big]e^{-\pi r^2t}\,t^2\,dt .
\]
\end{definition}

By Lemma~\ref{lem:psiasym} the second integral converges for every
$r^2$, and the first converges precisely when $r^2>2$ (the integrand
grows like $-728\,e^{(2-r^2)\pi t}t^2$).  Definition~\ref{def:Iplus}
is one natural regularization among several: any other subtraction
scheme differs from it by an \emph{explicit elementary function} of
$r^2$ (incomplete-gamma integrals of the subtracted terms).  Selecting
the canonical normalization --- the one for which the analytic
continuation below produces the linear-programming certificate --- is
part of Open Problem~\ref{op:continuation}, exactly as the
choice of contour is the heart of the dimension-$8$ proof
\cite[\S7]{Viazovska2017}.

\subsection{The candidate magic function}\label{subsec:f6}

\begin{definition}\label{def:f6}
For $\alpha,\beta\in\R$ and $x\in\R^6$ with $|x|^2>2$, set
\[
f_6(x)\;=\;\sin^2\!\Big(\frac{\pi|x|^2}{2}\Big)
\Big[\alpha\,\Ip(|x|^2)+\beta\,\Imm(|x|^2)\Big],
\]
extended to all of $\R^6$ by the conjectural analytic continuation of
$\Ip$ (Open Problem~\ref{op:continuation}), with $f_6(0)$ defined as
the limit $r\to0^+$.
\end{definition}

The design conditions on the pair $(\alpha,\beta)$ and the
continuation are:
\begin{equation}\label{eq:CE-design}
\begin{aligned}
\text{(CE1)}&\quad f_6 \text{ and } \fhat{f_6} \text{ continuous,
integrable, with rapid decay (admissibility);}\\
\text{(CE2)}&\quad f_6(x)\le0 \text{ for } |x|\ge\sqrt2
\text{ (at least on } \sqrt2\le|x|\le2\sqrt2\text{);}\\
\text{(CE3)}&\quad \fhat{f_6}(\xi)\ge0 \text{ for all } \xi;\\
\text{(CE4)}&\quad \fhat{f_6}(\sqrt2)=0;\\
\text{(CE5)}&\quad f_6(0)/\fhat{f_6}(0)=72 .
\end{aligned}
\end{equation}
Conditions (CE4)--(CE5) are two linear constraints on
$(\alpha,\beta)$ modulo overall scale, so they determine the candidate
projectively once the continuation is fixed.

\begin{conjecture}[Sharp certificate]\label{conj:certificate}
There is a canonical analytic continuation of $\Ip$ to $r^2\ge0$ and a
choice of $(\alpha,\beta)$ such that $f_6$ satisfies
\textup{(CE1)--(CE5)} and descends, via
Proposition~\ref{prop:certificate} (or via the interpolation framework
of \cite{CKMRV2022}), to a sharp certificate
\[
K(\R^6)\;\le\;72 .
\]
Together with the $E_6$ configuration this gives $K(\R^6)=72$.
\end{conjecture}

\subsection{Why each ingredient is forced}\label{subsec:forced}

We summarize why \eqref{eq:f6} is not one ansatz among many but the
output of a short forced search:

\begin{enumerate}
\item \emph{Level $3$.}  $\det E_6=3$ and
$\theta_{E_6}\in M_3(\Gamma_0(3),\chi_{-3})$
(Observation~\ref{obs:level}).
\item \emph{Weight $6$.}  The radial pairing in $\R^6$ uses the
measure $t^{(6-2)/2}dt=t^2dt$, which couples to weight-$6$
transformation laws under $t\mapsto1/(3t)$
(Lemma~\ref{lem:Fpositive} and Remark~\ref{rem:rescale}).
\item \emph{The cusp form.}  $\dim S_6(\Gamma_0(3))=1$: the form
$\Dthree=\eta^6\eta_3^6$ is the unique choice.
\item \emph{The $\pm$ decomposition.}  $W_3$ splits
$M_6(\Gamma_0(3))$ into eigenspaces (Lemma~\ref{lem:W3eisenstein});
the two transforms $\Ip,\Imm$ are modelled on the two Fourier
parities, as in dimension $8$.
\item \emph{The numerator.}  A Fricke-\emph{even} meromorphic
quotient with cusp-form denominator requires numerator eigenvalue
$-1$; among the Eisenstein bilinears $G_+^2,G_-^2,G_+G_-$ only
$G_+G_-$ qualifies (\S\ref{subsec:psiplus}).
\item \emph{The pole is necessary.}  A holomorphic even choice
(e.g.\ $G_+$ itself) yields a transform with no sign-changing
interpolation structure; in dimension $8$ it is precisely the simple
pole of the integrand at the cusp that, after contour deformation,
produces the interpolation basis and the sharp zeros
\cite{Viazovska2017,Cohn2017}.  The constant term of $G_+G_-$ at
$\infty$, namely $(1+27)(1-27)=-728\ne0$, guarantees the simple pole
(Proposition~\ref{prop:psiplus}).
\item \emph{The prefactor.}  $\sin^2(\pi r^2/2)$ vanishes to second
order exactly at $r^2\in2\Z_{>0}$.  The squared distances realized in
the $E_6$ kissing configuration are $|u-v|^2=4-2\langle u,v\rangle\in
\{2,4,6,8\}$, since distinct minimal vectors have
$\langle u,v\rangle\in\{1,0,-1,-2\}$: \emph{every} realized distance is a
zero of the prefactor, the complementary-slackness pattern of a sharp
certificate.  The expected singular terms of the continued bracket at
$r^2=0$ and $r^2=2$ (poles in $r^2$, in analogy with
\cite[\S7]{Viazovska2017}) convert the prefactor's double zeros at
those two radii into a positive value $f_6(0)>0$ and the prescribed
behavior at $r=\sqrt2$.
\end{enumerate}

%% ============================================================
\section{Verified properties of the building blocks}
\label{sec:verified}
%% ============================================================

This section records what is \emph{proved} and what is
\emph{computationally verified} about $\Dthree$ and $\psip$.  All
computations were performed twice, by independent code paths: first
with floating/symbolic series (mpmath/sympy), then re-derived from
scratch with exact integer arithmetic, and cross-validated against the
LMFDB \cite{LMFDB}.  The scripts and data files are ancillary to this
paper and are described in Appendix~\ref{app:repro}; we cite them by
name below.

\subsection{The cusp form is the newform \textsf{3.6.a.a}}

\begin{proposition}\label{prop:hecke}
$\Dthree=\eta(\tau)^6\eta(3\tau)^6$ spans $S_6(\Gamma_0(3))$ and is a
normalized Hecke eigenform; it is a newform (there are no oldforms,
since $S_6(\mathrm{SL}_2(\Z))=0$), self-dual, with trivial character.
Its Hecke eigenvalues $a_n$ (Appendix~\ref{app:hecke}) satisfy:
\begin{enumerate}
\item[(i)] multiplicativity $a_{mn}=a_ma_n$ for $(m,n)=1$, verified
exactly for all $160$ coprime pairs with $mn\le100$;
\item[(ii)] the weight-$6$ prime-power recursion
$a_{p^2}=a_p^2-p^5$, verified at $p\in\{2,5,7\}$, and the ramified
relation $a_9=a_3^2=81$ at the level prime;
\item[(iii)] a perfect match with the LMFDB newform \textsf{3.6.a.a}
for all listed coefficients $a_2,\dots,a_{31}$, including the derived
value $a_{25}=a_5^2-5^5=-3089$;
\item[(iv)] $a_3=+9=3^{(k-2)/2}$ with $k=6$, consistent with the
Atkin--Lehner relation of Proposition~\ref{prop:AL}.
\end{enumerate}
\end{proposition}

\begin{proof}
Modularity, weight, level, trivial character, and cuspidality of the
eta quotient follow from the standard criteria
\cite[\S1.4]{Ono2004}: for $\eta(\tau)^6\eta(3\tau)^6$ one has
$\sum_\delta \delta\,r_\delta=6+18=24\equiv0\ (24)$,
$\sum_\delta (N/\delta)\,r_\delta=18+6=24\equiv0\ (24)$, and
$\prod_\delta \delta^{r_\delta}=3^6$ is a perfect square, giving
trivial nebentypus; the cusp orders are positive.  Since
$\dim S_6(\Gamma_0(3))=1$ \cite{DiamondShurman2005,LMFDB} and the
Hecke operators preserve this line, $\Dthree$ is automatically an
eigenform, and $a_1=1$ normalizes it; multiplicativity and the
recursion are then theorems of Hecke theory, so the checks (i)--(ii)
validate the computed $q$-expansion rather than the theory.  The
verifications are performed with exact integer arithmetic in the
ancillary script \texttt{claudecode\_independent\_verify.py} (checks
C1--C3), independently of the original computation in
\texttt{compute\_hecke\_eigenvalues.py}; the LMFDB comparison (iii) is
recorded in \texttt{data/hecke\_eigenvalues.json}
\cite{LMFDB}.
\end{proof}

\subsection{The Atkin--Lehner sign}

\begin{proposition}\label{prop:AL}
$\Dthree|_6W_3=-\Dthree$; that is, the Atkin--Lehner (= Fricke)
eigenvalue of the newform \textsf{3.6.a.a} is $\eps_3=-1$.
\end{proposition}

\begin{proof}
From $\eta(-1/z)=\sqrt{-iz}\,\eta(z)$ applied at $z=3\tau$ and
$z=\tau$:
\[
\Dthree\!\big(-\tfrac1{3\tau}\big)
=\eta\!\big(-\tfrac1{3\tau}\big)^6\,\eta\!\big(-\tfrac1{\tau}\big)^6
=(-3i\tau)^3(-i\tau)^3\,\eta(3\tau)^6\eta(\tau)^6
=-27\,\tau^6\,\Dthree(\tau),
\]
since $(-3i\tau)^3(-i\tau)^3=(27i)(i)\tau^6=-27\tau^6$.  Hence
$(\Dthree|_6W_3)(\tau)=3^{-3}\tau^{-6}\cdot(-27\tau^6)\Dthree(\tau)
=-\Dthree(\tau)$.
\end{proof}

The sign was confirmed four independent ways: (a) the symbolic
eta-transformation derivation above, executed in sympy
(\texttt{verify\_atkin\_lehner.py}); (b) the Atkin--Lehner newform
relation $a_p=-\eps_p\,p^{k/2-1}$ for $p\,\|\,N$
\cite{AtkinLehner1970}, which at $p=3$, $k=6$ reads
$a_3=-9\,\eps_3$, so $a_3=+9$ forces $\eps_3=-1$; (c) a $30$-digit
numerical functional-equation test at six points on the imaginary
axis; (d) the LMFDB entry for \textsf{3.6.a.a} \cite{LMFDB}.  Items
(b)--(c) are checks C5--C6 of
\texttt{claudecode\_independent\_verify.py}.

\begin{remark}[Imaginary-axis sign conventions]\label{rem:signs}
A bookkeeping point that is easy to get wrong, recorded here because
it is load-bearing for any re-derivation.  With the normalization
\eqref{eq:slashW3}, evaluating at $\tau=it$ gives
$(it)^{-6}=-t^{-6}$, so for a weight-$6$ eigenfunction
$f|_6W_3=\eps f$ the \emph{observable} relation on the imaginary axis
is
\[
f\!\big(\tfrac{i}{3t}\big)\;=\;-\,\eps\cdot 27\,t^6 f(it).
\]
Thus the $\eps=-1$ form satisfies the \emph{positive} relation
$\Dthree(i/(3t))=+27t^6\Dthree(it)$ (consistent with
$\Dthree(it)>0$), while the $\eps=+1$ form satisfies
$\psip(i/(3t))=-27t^6\,\psip(it)$ --- which is what forces the zero in
Proposition~\ref{prop:forcedzero}.
\end{remark}

\subsection{Ramanujan--Petersson consistency}

\begin{proposition}\label{prop:RP}
The computed eigenvalues satisfy $|a_p|\le 2p^{5/2}$ for every prime
$p\le97$; the extremal ratio $|a_p|/2p^{5/2}$ on this range is
$0.825$, attained at $p=53$ (Appendix~\ref{app:hecke}).
\end{proposition}

The bound is of course a theorem for all $p$ (Deligne
\cite{Deligne1974}); the verification is a consistency test of the
computed expansion --- any coefficient error of even modest size would
violate it --- not new mathematics.  It is check C4 of
\texttt{claudecode\_independent\_verify.py}.

\subsection{Properties of \texorpdfstring{$\psip$}{psi+}}

\begin{proposition}\label{prop:psiplus}
The function $\psip$ of \eqref{eq:psiplus} is a meromorphic modular
form of weight $6$ on $\Gamma_0(3)$, holomorphic on $\HH$, with
\[
\psip|_6W_3\;=\;+\,\psip ,
\]
simple poles at both cusps, and integral Laurent coefficients.  At
$\infty$ its expansion is \eqref{eq:psilaurent}, with
\[
\res_{\infty}\psip\;:=\;[q^{-1}]\,\psip\;=\;-728\;=\;-2^3\cdot7\cdot13,
\qquad
[q^{0}]\,\psip\;=\;-5376\;=\;-2^8\cdot3\cdot7 ,
\]
and the expansion of $\psip$ in the local parameter at the cusp $0$
(i.e.\ of $\psip|_6W_3$ at $\infty$) is the same series; in particular
the pole at $0$ is also simple with leading coefficient $-728$.
\end{proposition}

\begin{proof}
Weight and level are clear from \eqref{eq:psiplus};
holomorphy on $\HH$ holds because $\eta$ is nonvanishing there.  The
eigenvalue: by Lemma~\ref{lem:W3eisenstein},
$(G_+G_-)|_{12}W_3=(G_+|_6W_3)(G_-|_6W_3)=-G_+G_-$, and by
Proposition~\ref{prop:AL} the denominator also picks up $-1$;
multiplicativity of the slash on products and quotients gives
$\psip|_6W_3=(-1)/(-1)\cdot\psip=\psip$.  Pole structure: $\Dthree$
has a simple zero at each cusp (\S\ref{subsec:blocks}) and the
numerator has constant term $(1+27)(1-27)=-728\ne0$ at $\infty$, so
the pole at $\infty$ is exactly simple; the eigenfunction property
transports the same Laurent series to the cusp $0$.  Integrality: the
numerator lies in $\Z[[q]]$ and
$\Dthree/q\in1+q\Z[[q]]$ has an inverse in $1+q\Z[[q]]$, so
$\psip\in q^{-1}\Z[[q]]$.  The expansion was computed by exact integer
series division to order $q^{80}$ (all coefficients integral, a
consistency check on the division) in
\texttt{claudecode\_independent\_verify.py} (checks C7--C8), with the
numerator independently verified in
\texttt{verify\_psi\_plus\_residue.py}; the eigenvalue was also
confirmed numerically at six imaginary-axis points at $30$ digits
(check C9).
\end{proof}

\subsection{A forced zero at the Fricke fixed point}

The fixed point of $W_3$ in $\HH$ is $\tau_*=i/\sqrt3$
(indeed $-1/(3\tau_*)=\tau_*$).

\begin{proposition}\label{prop:forcedzero}
Any weight-$6$ meromorphic form $f$ on $\Gamma_0(3)$ with
$f|_6W_3=+f$ that is finite at $\tau_*=i/\sqrt3$ vanishes there.  In
particular
\[
\psip\!\Big(\frac{i}{\sqrt3}\Big)\;=\;0 ,
\]
while $\Dthree(i/\sqrt3)=0.02251\ldots\ne0$ (the $\eps=-1$ form has no
forced zero).  Numerically,
$|\psip(i/\sqrt3)|/|\psip(0.75\,i)|\approx4\times10^{-32}$ at
$30$-digit working precision, confirming the identity to the full
precision of the computation.
\end{proposition}

\begin{proof}
Evaluate the eigenfunction property at the fixed point:
$f(\tau_*)=(f|_6W_3)(\tau_*)=3^{-3}\tau_*^{-6}f(\tau_*)$, and
$\tau_*^6=(i/\sqrt3)^6=i^6/27=-1/27$, so $3^{-3}\tau_*^{-6}=-1$ and
$f(\tau_*)=-f(\tau_*)$.  (Equivalently: set $t=1/\sqrt3$ in
$\psip(i/(3t))=-27t^6\psip(it)$ of Remark~\ref{rem:signs}.)
$\psip$ is holomorphic at $\tau_*\in\HH$ by
Proposition~\ref{prop:psiplus}.  The numerical confirmation is check
C10 of \texttt{claudecode\_independent\_verify.py}.
\end{proof}

\begin{remark}\label{rem:zerorelevance}
By Lemma~\ref{lem:psiasym}, $\psip(it)$ is positive near $t=0$ and
negative for large $t$; Proposition~\ref{prop:forcedzero} pins a sign
change at exactly $t_*=1/\sqrt3$, the fixed point of
$t\mapsto1/(3t)$.  This is directly relevant to the open continuation
problem: in dimension $8$ the deformed contour pivots at the fixed
point $\tau=i$ of the inversion \cite[\S7]{Viazovska2017}, and the
natural level-$3$ contour pivots at $\tau_*=i/\sqrt3$ --- where, by
the proposition, the integrand \emph{vanishes}.  We expect this forced
zero to serve as the anchor of the positivity analysis of the
continued $\Ip$ (Open Problem~\ref{op:continuation}).
\end{remark}

\subsection{Arithmetic of the Laurent coefficients}
\label{subsec:laurentarith}

\begin{proposition}\label{prop:fiftysix}
Write $\psip=\sum_{n\ge-1}c_n q^n$.  Then $56\mid c_n$ for every
$n\ge-1$.
\end{proposition}

\begin{proof}
Let $N(q)=E_6(\tau)^2-729E_6(3\tau)^2=\sum_{n\ge0}N_nq^n$.  Since
$E_6\equiv1\pmod{504}$ coefficientwise ($504=2^3\cdot3^2\cdot7$),
every coefficient of $E_6(\tau)^2$ and of $E_6(3\tau)^2$ beyond the
constant term is divisible by $1008=2\cdot504$
(cross terms $2\cdot(-504)\sigma_5$ and squares $504^2(\cdots)$), so
$N_0=1-729=-728$ and $1008\mid N_n$ for $n\ge1$.  Writing
$q\,\psip=N\cdot U$ with
$U=(\Dthree/q)^{-1}\in1+q\Z[[q]]$, every coefficient of $q\,\psip$ is
$-728\,u_m+1008\,(\text{integer})$, and
$\gcd(728,1008)=56$ divides both terms.
\end{proof}

The divisibility was confirmed empirically for all computed
coefficients $c_{-1},\dots,c_{80}$.  Note $56=2^3\cdot7$.

\begin{observation}\label{obs:alphabet}
The two distinguished coefficients of $\psip$ --- the residue and the
constant term, which are exactly the data entering the cusp residue
calculus of the continuation problem --- factor over the prime set
$\{2,3,5,7,11,13\}$ of Observation~\ref{obs:smooth}:
\[
c_{-1}=-728=-2^3\cdot7\cdot13,
\qquad
c_0=-5376=-2^8\cdot3\cdot7 .
\]
The regular coefficients do not: already
$c_1=195048=2^3\cdot3^4\cdot7\cdot43$ contains the prime $43$
(Appendix~\ref{app:laurent}).  We state this honestly: the small-prime
structure governs precisely the principal part (plus constant term)
--- the analytically load-bearing data --- and, by
Proposition~\ref{prop:fiftysix}, a fixed factor $56=2^3\cdot7$ of
every coefficient; it does not govern the full $q$-expansion.  The
regular coefficients are not multiplicative in any tested sense
(e.g.\ $c_2c_3\ne c_6$), as expected for a meromorphic form that is
not a Hecke eigenform.
\end{observation}

\subsection{Numerical profile of the cusp-form transform}
\label{subsec:Iminusnum}

$\Imm(r^2)$ was computed by adaptive quadrature at $30$-digit working
precision at $r^2\in\{0,1,\dots,8,10,12,15,20\}$
(\texttt{compute\_I\_minus.py}; values in
Appendix~\ref{app:Iminus}).  The profile is smooth, positive and
monotonically decreasing, with $\Imm(0)\approx4.5155\times10^{-3}$,
and the tail decays at an exponential-type rate consistent with the
saddle-point prediction
$\log \Imm(r^2)\sim-\tfrac{2\sqrt6}{3}\pi r$ (and with the rigorous
bound of Lemma~\ref{lem:Iminus}); for instance
$\Imm(20)\approx4.17\times10^{-10}$.  At even integer values
$r^2\in\{2,4,6,8,\dots\}$ --- the squared distances realized by the
$E_6$ configuration --- the prefactor $\sin^2(\pi r^2/2)$ of
\eqref{eq:f6} vanishes to second order, so the assembled candidate
automatically interpolates zero at every realized distance, the
complementary-slackness pattern of \S\ref{subsec:forced}(7).

%% ============================================================
\section{The analytic continuation gap}\label{sec:gap}
%% ============================================================

We now state precisely what separates
Conjecture~\ref{conj:certificate} from a theorem.  \emph{None of the
following is established in this paper.}

\begin{openproblem}[Canonical continuation]\label{op:continuation}
Construct the analytic continuation of $\Ip(r^2)$ from $r^2>2$ to all
$r^2\ge0$, together with its canonical normalization.  Concretely:
deform the ray $(0,i\infty)$ to a contour pivoting at the Fricke fixed
point $\tau_*=i/\sqrt3$ (where the integrand vanishes,
Proposition~\ref{prop:forcedzero}), and express the continued
$\Ip$ as a deformed integral plus the explicit contributions of the
principal parts $-728\,q^{-1}-5376$ at the two cusps of
$\Gamma_0(3)$.  By analogy with \cite[\S7]{Viazovska2017} the
continuation is expected to carry singular elementary terms at
$r^2=0$ and $r^2=2$ (poles in $r^2$) which cancel the double zeros of
the prefactor $\sin^2(\pi r^2/2)$ at exactly those two radii, leaving
$f_6(0)>0$ and the prescribed sharp behavior at $r=\sqrt2$.  The
residue data are integral and highly structured
(Proposition~\ref{prop:fiftysix}, Observation~\ref{obs:alphabet});
the two-cusp geometry and the $\sqrt3$-rescaling of
Remark~\ref{rem:rescale} are the genuinely new features relative to
level $1$.
\end{openproblem}

\begin{openproblem}[Regularity and decay]\label{op:schwartz}
Show that the continued bracket
$B(r^2)=\alpha\Ip(r^2)+\beta\Imm(r^2)$ is smooth on $(0,\infty)$,
that $f_6$ extends continuously to $x=0$, and that $f_6$ and
$\fhat{f_6}$ satisfy the admissibility/decay condition (CE1).  (The
cusp-form component is unproblematic: Lemma~\ref{lem:Iminus} already
gives exponential-type decay; the burden falls on the continued
$\Ip$.)
\end{openproblem}

\begin{openproblem}[Sign conditions]\label{op:signs}
Verify (CE2) and (CE3): $f_6\le0$ on $|x|\ge\sqrt2$ (at minimum on
the kissing annulus $\sqrt2\le|x|\le2\sqrt2$) and
$\fhat{f_6}\ge0$ everywhere.  In dimension $8$ the analogous step
required delicate interval-arithmetic and ad hoc inequalities; the
anchor here should be the forced zero at the contour pivot
(Remark~\ref{rem:zerorelevance}) and the strict positivity of the
cusp-form component (Lemma~\ref{lem:Fpositive}).
\end{openproblem}

\begin{openproblem}[Normalization]\label{op:alphabeta}
Determine $(\alpha,\beta)$ from the linear system (CE4)--(CE5), in
closed form.  In dimension $8$ the analogous constants are explicit
elementary multiples of powers of $\pi$; we conjecture the same here,
with the integers $728$, $5376$, and $27$ entering through the residue
calculus.
\end{openproblem}

\begin{openproblem}[Descent to the sphere]\label{op:descent}
Convert the Euclidean function $f_6$ into a kissing bound: either
compute $c_0(f_6)>0$ and apply Proposition~\ref{prop:certificate},
verifying $f_6(0)/c_0(f_6)=72$ exactly, or run the candidate through
an interpolation argument in the style of \cite{CKMRV2022}.  This step
is where the design ratio (CE5) must be reconciled with the spherical
average $c_0$, and we flag it as a genuine (if presumably lighter)
part of the program rather than a formality.
\end{openproblem}

\subsection*{What we do not claim}

For the record, and in deliberate redundancy with
\S\ref{subsec:noclaim}:

\begin{itemize}
\item We do \emph{not} claim a proof of $K(\R^6)=72$.  We claim an
explicit candidate plus verified building blocks plus a precise
statement of the remaining analytic problem.
\item Numerical evidence --- including any future sampling of $f_6$
at finitely many radii --- cannot settle
Open Problems~\ref{op:signs}--\ref{op:descent}; sign conditions must
hold at \emph{all} radii, and only the continuation argument can
deliver that.
\item The structural observations of Section~\ref{sec:structural}
motivated the construction but carry no evidential weight for the
truth of Conjecture~\ref{conj:main}.  The candidate stands or falls
with the analysis of Open
Problems~\ref{op:continuation}--\ref{op:descent}.
\end{itemize}

%% ============================================================
\section{Comparison with the dimension-eight construction}
\label{sec:comparison}
%% ============================================================

Table~\ref{tab:compare} aligns the present construction with
Viazovska's.  The columns are deliberately parallel: every level-$1$
ingredient has a unique level-$3$ counterpart, and the two genuinely
new difficulties (two cusps instead of one; the $\sqrt3$-rescaling
between a lattice and its dual) are confined to the last two rows.

\begin{table}[ht]
\centering
\small
\begin{tabular}{>{\raggedright\arraybackslash}p{0.30\linewidth}
                >{\raggedright\arraybackslash}p{0.32\linewidth}
                >{\raggedright\arraybackslash}p{0.30\linewidth}}
\toprule
 & \textbf{Dimension 8} (theorem) & \textbf{Dimension 6} (this paper,
conjectural) \\
\midrule
Lattice & $E_8$, unimodular & $E_6$, $\det=3$ \\
Minimal vectors & $240$ & $72$ \\
Dual configuration & self-dual & $E_6^*\supset E_6$, index $3$;
$54$ minimal vectors \\
Modular group & $\mathrm{SL}_2(\Z)$ & $\Gamma_0(3)$ \\
Cusps & one ($\infty$) & two ($0$, $\infty$) \\
Theta series & $\theta_{E_8}=E_4$ & $\theta_{E_6}\in
M_3(\Gamma_0(3),\chi_{-3})$ \\
Distinguished cusp form & $\Delta=\eta^{24}$, weight $12$,
$\dim S_{12}=1$ & $\Dthree=\eta^6\eta_3^6$, weight $6$,
$\dim S_6(\Gamma_0(3))=1$ \\
Eisenstein input & $E_2,E_4,E_6$ (level $1$) &
$G_\pm=E_6(\tau)\pm27E_6(3\tau)$ \\
Meromorphic seed & quotients with $\Delta$ in the denominator &
$\psip=G_+G_-/\Dthree$ \\
Parity mechanism & inversion $S:\tau\mapsto-1/\tau$ &
Fricke $W_3:\tau\mapsto-1/(3\tau)$ \\
Fixed point of involution & $\tau=i$ & $\tau_*=i/\sqrt3$;
$\psip(\tau_*)=0$ forced \\
Pole data driving residues & $q$-expansion principal parts at the
single cusp & $-728\,q^{-1}-5376$ at \emph{each} of two cusps \\
Transform & Laplace transforms of modular seeds; contour pivots at
$i$ & $t^2\,dt$ transforms (Definitions~\ref{def:Iminus},
\ref{def:Iplus}); contour expected to pivot at $i/\sqrt3$ \\
Self-duality of transform & Fricke parity $=$ Fourier parity &
parity twisted by a $\sqrt3$-rescaling
(Remark~\ref{rem:rescale}) \\
Status & sphere packing solved \cite{Viazovska2017}; kissing $240$
classical \cite{Levenshtein1979,OdlyzkoSloane1979}; universal
optimality \cite{CKMRV2022} & kissing $72$ conjectured; building
blocks verified (\S\ref{sec:verified}); continuation open
(\S\ref{sec:gap}) \\
\bottomrule
\end{tabular}
\medskip
\caption{The level-$3$ construction as the structural analogue of the
level-$1$ construction.}
\label{tab:compare}
\end{table}

Two remarks on the comparison.  First, in dimension $8$ the kissing
number itself never needed modular forms: a degree-six polynomial
certificate already matches $240$
\cite{Levenshtein1979,OdlyzkoSloane1979}.  Dimension $6$ is harder in
precisely this sense: polynomial and SDP certificates provably stall
at $77$--$78$, so if $72$ is the truth, the certificate must be
transcendental --- which is why a Viazovska-type construction is the
natural (perhaps the only known) escalation.  Second, the level-$3$
problem is not a routine transcription.  The two-cusp residue
bookkeeping, the $\eps=\pm1$ eigenspace split across a
\emph{two}-dimensional Eisenstein space, and the dual-lattice
rescaling in Remark~\ref{rem:rescale} have no level-$1$ counterparts;
they are the concrete content of
Open Problem~\ref{op:continuation}.

%% ============================================================
\section{Conclusion}\label{sec:conclusion}
%% ============================================================

We have proposed an explicit candidate magic function for the
six-dimensional kissing problem: a $\sin^2$-weighted combination of
$t^2dt$-Laplace transforms of the unique weight-$6$ cusp form
$\Dthree=\eta^6\eta_3^6$ on $\Gamma_0(3)$ and of the forced
Fricke-even meromorphic quotient
$\psip=(E_6^2-729E_6(3\cdot)^2)/\Dthree$.  Every discrete choice in
the construction is dictated by the arithmetic of $E_6$ (level,
weight, cusp form, eigenspace, numerator), and every property the
construction demands of its building blocks has been proved or
verified by independent computations cross-checked against the LMFDB,
including two facts we have not found elsewhere: the forced vanishing
$\psip(i/\sqrt3)=0$ at the Fricke fixed point, and the uniform
divisibility $56\mid c_n$ of the Laurent coefficients.  What remains
is the analytic core --- the canonical continuation of $\Ip$ with its
two-cusp residue calculus, the sign conditions, and the
normalization --- which we have isolated as
Open Problems~\ref{op:continuation}--\ref{op:descent}.  We believe
this is the correct division of labor for a problem of this type: the
discrete search is finished and documented; the continuation argument
can now be attacked in isolation.  The same template (discriminant
$d$ $\Rightarrow$ level $d$; dimension $n$ $\Rightarrow$ weight and
measure; one-dimensional cusp space $\Rightarrow$ forced seed) is
available in other open dimensions, and we intend to return to the
continuation problem in future work.

\subsection*{Data and code availability}
All verification scripts and data files accompany this paper as
ancillary files and are described in Appendix~\ref{app:repro}.  The
modular-form data were cross-validated against the LMFDB
\cite{LMFDB} (newform \textsf{3.6.a.a}, accessed June 10, 2026).

%% TODO: optional figures from the visualization session
%% (../figures/dim6_magic_function_shape.png,
%%  ../figures/eta_cusp_form_shape.png,
%%  ../figures/f6_real_space_shape.png) --- enable if the venue wants
%% illustrative figures; paths must be adjusted to the submission
%% directory layout.
%% \begin{figure}[ht]
%% \centering
%% \includegraphics[width=0.9\linewidth]{eta_cusp_form_shape.png}
%% \caption{$\Dthree$ on the upper half-plane.}
%% \end{figure}

%% ============================================================
\appendix
%% ============================================================

\section{Hecke eigenvalues of
\texorpdfstring{$\Dthree=\eta^6\eta_3^6$}{Delta3} for
\texorpdfstring{$p\le97$}{p<=97}}\label{app:hecke}

Exact values from integer-arithmetic expansion
(\texttt{claudecode\_independent\_verify.py}); entries
$a_2,\dots,a_{31}$ match the LMFDB newform \textsf{3.6.a.a}
\cite{LMFDB} exactly.  The last column is the Ramanujan--Petersson
ratio $|a_p|/2p^{5/2}<1$ (Proposition~\ref{prop:RP}).

\begin{center}
\small
\begin{tabular}{rrlc}
\toprule
$p$ & $a_p$ & factorization & $|a_p|/2p^{5/2}$ \\
\midrule
$2$ & $-6$ & $-2\cdot3$ & $0.530$ \\
$3$ & $9$ & $3^2$ & $0.289$ \\
$5$ & $6$ & $2\cdot3$ & $0.054$ \\
$7$ & $-40$ & $-2^3\cdot5$ & $0.154$ \\
$11$ & $-564$ & $-2^2\cdot3\cdot47$ & $0.703$ \\
$13$ & $638$ & $2\cdot11\cdot29$ & $0.524$ \\
$17$ & $882$ & $2\cdot3^2\cdot7^2$ & $0.370$ \\
$19$ & $-556$ & $-2^2\cdot139$ & $0.177$ \\
$23$ & $-840$ & $-2^3\cdot3\cdot5\cdot7$ & $0.166$ \\
$29$ & $4638$ & $2\cdot3\cdot773$ & $0.512$ \\
$31$ & $4400$ & $2^4\cdot5^2\cdot11$ & $0.411$ \\
$37$ & $-2410$ & $-2\cdot5\cdot241$ & $0.145$ \\
$41$ & $-6870$ & $-2\cdot3\cdot5\cdot229$ & $0.319$ \\
$43$ & $9644$ & $2^2\cdot2411$ & $0.398$ \\
$47$ & $-18672$ & $-2^4\cdot3\cdot389$ & $0.617$ \\
$53$ & $33750$ & $2\cdot3^3\cdot5^4$ & $0.825$ \\
$59$ & $-18084$ & $-2^2\cdot3\cdot11\cdot137$ & $0.338$ \\
$61$ & $39758$ & $2\cdot103\cdot193$ & $0.684$ \\
$67$ & $-23068$ & $-2^2\cdot73\cdot79$ & $0.314$ \\
$71$ & $-4248$ & $-2^3\cdot3^2\cdot59$ & $0.050$ \\
$73$ & $-41110$ & $-2\cdot5\cdot4111$ & $0.452$ \\
$79$ & $21920$ & $2^5\cdot5\cdot137$ & $0.198$ \\
$83$ & $82452$ & $2^2\cdot3\cdot6871$ & $0.657$ \\
$89$ & $-94086$ & $-2\cdot3^2\cdot5227$ & $0.630$ \\
$97$ & $49442$ & $2\cdot59\cdot419$ & $0.267$ \\
\bottomrule
\end{tabular}
\end{center}

\begin{remark}
A curiosity we record without any claim: the prime factors of $a_p$
for $p\le17$ lie in Ogg's set of fifteen supersingular primes
$\{2,3,5,7,11,13,17,19,23,29,31,41,47,59,71\}$ (first exit at
$a_{19}=-2^2\cdot139$).  The eta quotient $\eta(\tau)^6\eta(3\tau)^6$
is one of the classical multiplicative eta quotients
\TODO{confirm membership and numbering in Martin's table
\cite{Martin1996}}.
\end{remark}

\section{Laurent coefficients of \texorpdfstring{$\psip$}{psi+}}
\label{app:laurent}

Exact integer values from series division
(\texttt{claudecode\_independent\_verify.py}, check C8; full table to
$q^{80}$ in \texttt{data/claudecode\_independent.json}).  Every entry
is divisible by $56$ (Proposition~\ref{prop:fiftysix}).

\begin{center}
\small
\begin{tabular}{rrl}
\toprule
$n$ & $c_n$ & factorization \\
\midrule
$-1$ & $-728$ & $-2^3\cdot7\cdot13$ \\
$0$ & $-5376$ & $-2^8\cdot3\cdot7$ \\
$1$ & $195048$ & $2^3\cdot3^4\cdot7\cdot43$ \\
$2$ & $18475520$ & $2^9\cdot5\cdot7\cdot1031$ \\
$3$ & $508641336$ & $2^3\cdot3^3\cdot7\cdot336403$ \\
$4$ & $7509293568$ & $2^9\cdot3^4\cdot7\cdot25867$ \\
$5$ & $74665421768$ & $2^3\cdot7\cdot41\cdot32519783$ \\
$6$ & $565780723200$ & $2^9\cdot3^4\cdot5^2\cdot7\cdot11\cdot19\cdot373$ \\
$7$ & $3505407582096$ & $2^4\cdot3^4\cdot7\cdot71\cdot5442233$ \\
$8$ & $18588498442752$ & $2^9\cdot3\cdot7\cdot18341\cdot94261$ \\
$9$ & $87180602357016$ & $2^3\cdot3^3\cdot7\cdot47\cdot233\cdot5265193$ \\
$10$ & $369668580793344$ & $2^{10}\cdot3^6\cdot7\cdot421\cdot168037$ \\
\bottomrule
\end{tabular}
\end{center}

\section{Numerical values of \texorpdfstring{$\Imm(r^2)$}{I-(r2)}}
\label{app:Iminus}

Adaptive quadrature at $30$-digit working precision
(\texttt{compute\_I\_minus.py}; \texttt{data/I\_minus\_values.json});
$12$ significant digits shown.

\begin{center}
\small
\begin{tabular}{rl}
\toprule
$r^2$ & $\Imm(r^2)$ \\
\midrule
$0$ & $4.51552678207\times10^{-3}$ \\
$1$ & $8.17231376350\times10^{-4}$ \\
$2$ & $2.06065543421\times10^{-4}$ \\
$3$ & $6.29431588663\times10^{-5}$ \\
$4$ & $2.18574953063\times10^{-5}$ \\
$5$ & $8.33305293842\times10^{-6}$ \\
$6$ & $3.41354566353\times10^{-6}$ \\
$7$ & $1.48103225202\times10^{-6}$ \\
$8$ & $6.73715632719\times10^{-7}$ \\
$10$ & $1.56230062439\times10^{-7}$ \\
$12$ & $4.08403374890\times10^{-8}$ \\
$15$ & $6.46880116651\times10^{-9}$ \\
$20$ & $4.16885311026\times10^{-10}$ \\
\bottomrule
\end{tabular}
\end{center}

\section{Reproducibility}\label{app:repro}

All claims labelled ``verified computationally'' are reproduced by the
following ancillary files (Python~3; dependencies: \texttt{sympy},
\texttt{mpmath}).  Two independent code paths were used: the original
derivation scripts, and an independent re-derivation in exact integer
arithmetic.

\begin{itemize}
\item \texttt{verify\_atkin\_lehner.py} --- symbolic eta-transformation
derivation of $\Dthree|_6W_3=-\Dthree$ (Proposition~\ref{prop:AL}(a)).
\item \texttt{compute\_hecke\_eigenvalues.py} --- $q$-expansion of
$\Dthree$, eigenvalue table, multiplicativity and
Ramanujan--Petersson checks; output in
\texttt{data/hecke\_eigenvalues.json}.
\item \texttt{verify\_psi\_plus\_residue.py} --- expansion of the
numerator $G_+G_-$; residue $-728$.
\item \texttt{compute\_I\_minus.py} --- $30$-digit quadrature of
$\Imm(r^2)$; output in \texttt{data/I\_minus\_values.json}.
\item \texttt{claudecode\_independent\_verify.py} --- independent
exact-integer re-derivation; ten checks (C1 expansion sanity; C2
multiplicativity at $160$ coprime pairs; C3 prime-power recursion; C4
Ramanujan--Petersson; C5 Atkin--Lehner via the newform relation; C6
Atkin--Lehner via the numerical functional equation; C7 residue; C8
Laurent expansion of $\psip$ to $q^{80}$ with integrality; C9 Fricke
$+1$ functional equation of $\psip$ at six points; C10 forced zero at
$i/\sqrt3$), all passing; output in
\texttt{data/claudecode\_independent.json}.
\end{itemize}

External database: LMFDB newform \textsf{3.6.a.a},
\url{https://www.lmfdb.org/ModularForm/GL2/Q/holomorphic/3/6/a/a/},
accessed June 10, 2026 \cite{LMFDB}.

%% ============================================================
%% Bibliography
%% ============================================================

\begin{thebibliography}{99}

\bibitem{AtkinLehner1970}
A.~O.~L. Atkin and J. Lehner,
\emph{Hecke operators on $\Gamma_0(m)$},
Math. Ann. \textbf{185} (1970), 134--160.

\bibitem{BachocVallentin2008}
C. Bachoc and F. Vallentin,
\emph{New upper bounds for kissing numbers from semidefinite
programming},
J. Amer. Math. Soc. \textbf{21} (2008), 909--924.

\bibitem{BannaiSloane1981}
E. Bannai and N.~J.~A. Sloane,
\emph{Uniqueness of certain spherical codes},
Canad. J. Math. \textbf{33} (1981), 437--449.

\bibitem{CohnElkies2003}
H. Cohn and N. Elkies,
\emph{New upper bounds on sphere packings I},
Ann. of Math. (2) \textbf{157} (2003), 689--714.

\bibitem{Cohn2017}
H. Cohn,
\emph{A conceptual breakthrough in sphere packing},
Notices Amer. Math. Soc. \textbf{64} (2017), 102--115.

\bibitem{CKMRV2017}
H. Cohn, A. Kumar, S.~D. Miller, D. Radchenko, and M. Viazovska,
\emph{The sphere packing problem in dimension 24},
Ann. of Math. (2) \textbf{185} (2017), 1017--1033.

\bibitem{CKMRV2022}
H. Cohn, A. Kumar, S.~D. Miller, D. Radchenko, and M. Viazovska,
\emph{Universal optimality of the $E_8$ and Leech lattices and
interpolation formulas},
Ann. of Math. (2) \textbf{196} (2022), 983--1082.

\bibitem{ConwaySloane1999}
J.~H. Conway and N.~J.~A. Sloane,
\emph{Sphere Packings, Lattices and Groups},
3rd ed., Grundlehren der mathematischen Wissenschaften 290,
Springer-Verlag, New York, 1999.

\bibitem{Deligne1974}
P. Deligne,
\emph{La conjecture de Weil. I},
Inst. Hautes \'Etudes Sci. Publ. Math. \textbf{43} (1974), 273--307.

\bibitem{DGS1977}
P. Delsarte, J.~M. Goethals, and J.~J. Seidel,
\emph{Spherical codes and designs},
Geom. Dedicata \textbf{6} (1977), 363--388.

\bibitem{DiamondShurman2005}
F. Diamond and J. Shurman,
\emph{A First Course in Modular Forms},
Graduate Texts in Mathematics 228, Springer-Verlag, New York, 2005.

\bibitem{Ebeling2013}
W. Ebeling,
\emph{Lattices and Codes: A Course Partially Based on Lectures by
Friedrich Hirzebruch},
3rd ed., Springer Spektrum, Wiesbaden, 2013.

\bibitem{KabatianskyLevenshtein1978}
G.~A. Kabatiansky and V.~I. Levenshtein,
\emph{Bounds for packings on the sphere and in space},
Problemy Peredachi Informatsii \textbf{14} (1978), no. 1, 3--25.

\bibitem{Levenshtein1979}
V.~I. Levenshtein,
\emph{On bounds for packings in $n$-dimensional Euclidean space},
Dokl. Akad. Nauk SSSR \textbf{245} (1979), 1299--1303;
English transl., Soviet Math. Dokl. \textbf{20} (1979), 417--421.

\bibitem{LMFDB}
The LMFDB Collaboration,
\emph{The L-functions and Modular Forms Database},
\url{https://www.lmfdb.org}, 2026.
[Newform \textsf{3.6.a.a}; online, accessed June 10, 2026.]

\bibitem{MachadoOliveira2018}
F.~C. Machado and F.~M. de Oliveira Filho,
\emph{Improving the semidefinite programming bound for the kissing
number by exploiting polynomial symmetry},
Exp. Math. \textbf{27} (2018), 362--369.
\TODO{confirm volume/pages and the exact dimension-6 bound reported}

\bibitem{Martin1996}
Y. Martin,
\emph{Multiplicative $\eta$-quotients},
Trans. Amer. Math. Soc. \textbf{348} (1996), 4825--4856.
\TODO{confirm that $\eta(\tau)^6\eta(3\tau)^6$ appears in Martin's
classification table}

\bibitem{MittelmannVallentin2010}
H.~D. Mittelmann and F. Vallentin,
\emph{High-accuracy semidefinite programming bounds for kissing
numbers},
Exp. Math. \textbf{19} (2010), 175--179.
\TODO{confirm volume/pages and the dimension-6 value used in
\S\ref{subsec:dimsix}}

\bibitem{Musin2008}
O.~R. Musin,
\emph{The kissing number in four dimensions},
Ann. of Math. (2) \textbf{168} (2008), 1--32.

\bibitem{OdlyzkoSloane1979}
A.~M. Odlyzko and N.~J.~A. Sloane,
\emph{New bounds on the number of unit spheres that can touch a unit
sphere in $n$ dimensions},
J. Combin. Theory Ser. A \textbf{26} (1979), 210--214.

\bibitem{Ono2004}
K. Ono,
\emph{The Web of Modularity: Arithmetic of the Coefficients of Modular
Forms and $q$-series},
CBMS Regional Conference Series in Mathematics 102, Amer. Math. Soc.,
Providence, RI, 2004.

\bibitem{Schoenberg1942}
I.~J. Schoenberg,
\emph{Positive definite functions on spheres},
Duke Math. J. \textbf{9} (1942), 96--108.

\bibitem{SchuetteWaerden1953}
K. Sch\"utte and B.~L. van der Waerden,
\emph{Das Problem der dreizehn Kugeln},
Math. Ann. \textbf{125} (1953), 325--334.

\bibitem{Viazovska2017}
M.~S. Viazovska,
\emph{The sphere packing problem in dimension 8},
Ann. of Math. (2) \textbf{185} (2017), 991--1015.

\end{thebibliography}

\end{document}
